\chapter{Das heutige Theater}
\quellseite{2}

Viel wird gegenwärtig von dem erzieherischen und bildenden Einfluß des Theaters geredet. Der bekannte New Yorker Schauspieldirector Conried hat voriges Jahr öffentlich gesagt, das Theater könnte und sollte die Stelle der Kirche als Lehrerin des Volkes einnehmen; denn es seien allabendlich viel mehr Menschen in den Theatern New Yorks versammelt, als Sonntags in allen Kirchen der Stadt. Verschiedene amerikanische Prediger, namentlich der Episkopalkirche, haben mit dem Theater geradezu einen Bund geschlossen und wollen mit ihm zusammenarbeiten zur Hebung und Veredelung der Menschheit; andere Sectenprediger haben in letzter Zeit in ihrer Kirche ein Theater eingerichtet und hoffen, dadurch mehr an das Volk heranzukommen als durch ihre Predigt. Wie steht es nun mit dieser Volkserziehung, die das Theater bringen soll?

Wir wollen diesmal nicht selbst weiter ausführen, warum wir das heutige Theater für ganz verderblich und höchst verderblich halten. Wir wollen vielmehr andere Leute, die die heutige Bühne genau kennen, zu Worte kommen lassen. Wenn wir nämlich so scharf über das heutige Theater urtheilen und den Besuch desselben als ein Stück sündliches Weltwesen bezeichnen, an dem Christen keinen Theil haben sollen, wird gewöhnlich entgegnet, daß es doch auch unanständige Theaterstücke gebe. Ja, es gibt unanständige Stücke; aber solche Stücke werden selten oder gar nicht aufgeführt, und wenn sie aufgeführt werden, so werden sie als langweilig nicht besucht. Was denn heutigen Theater sein Gepräge gibt, das sind eben die anstößigen, unsittlichen, schändlichen Stücke, die von den Volksmassen angesehen und angehört werden. Das weiß jeder, der in einer Stadt lebt, die Theaterzettel sieht und die Besprechungen der Stücke in der täglichen Presse liest. Das bestätigen auch die Aussprachen der beiden Männer, die wir zu Worte kommen lassen wollen.

Ein bekannter New Yorker Prediger und Schriftsteller, Cyrus Townsend Brady, wollte im vorigen Winter einmal das heutige Theater und seinen „bildenden Einfluß“ aus eigener Anschauung kennen lernen. Er besuchte deshalb — daß ein Prediger freilich auch in solchem Interesse das Theater nicht besuchen sollte, ist gewiß — 21 Vorstellungen, und zwar nur in den besten Theatern, Komische Opern, sogenannte Zauberbildervorstellungen und vollends die auch bei der Welt als schlechte Theater geltenden Stätten 2c. hat er überhaupt nicht in seine Untersuchung gezogen, sonst wäre das Ergebniß noch ein ganz anderes geworden. Was Brady in den besten Theatern gefunden hat, hat er dann vor einigen Monaten in einer bekannten Zeitschrift näher geschildert.

Von den 21 Stücken, die er sich angesehen und angehört hat, bezeichnet er acht als unanständig. Uns ist freilich sehr zweifelhaft, ob alle acht wirklich durchaus unanständig waren, und ob nicht auch in dem einen oder andern dieser Stücke Sachen vorkamen, die ein ernstchristliches Gewissen verwerfen muß. Wir wissen von Leuten, die das heutige Theater genauer kennen, daß sie keine drei völlig unanständigen Stücke nennen können. Aber hören wir, was Brady weiter sagt: Vier dieser acht unanständigen Stücke, also gerade die Hälfte, waren so tödtlich langweilig, daß selbst das gute Spiel der Schauspieler ihnen kein Interesse verleihen konnte. Sie wirkten einschläfernd. Wer geht aber wohl in das Theater, um gelangweilt zu werden? Unter den vier andern unanständigen Stücken befanden sich Shakespeares „Julius Cäsar“, zwei Kriegsstücke und ein Drama, das in dem mittelalterlichen Rom spielte.

Sodann waren elf Theaterstücke, die nach Bradys Urtheil in der einen oder anderen Hinsicht durchaus anstößig waren. In jedem derselben spielte unerlaubte Liebe, also Unzucht, eine Hauptrolle, und Sachen wurden auf der Bühne gesagt oder angedeutet, die man in keiner weltlich ehrbaren Gesellschaft wiederholen könnte. Von diesen elf Stücken waren nur zwei uninteressant, nur eins wurde schlecht aufgeführt. In dem allerschlimmsten Stück spielte die beste Schauspielerin mit der besten Schauspielertruppe. Das ist ja gewöhnlich so. Gerade deshalb hat das heutige Theater eine solche verhängnißvolle Wirkung, weil die berühmtesten Schauspieler in den anstößigsten Stücken auftreten. Zwei der elf in Betracht kommenden Stücke waren große Opern, die also auch unter das Verwerfungsurtheil fallen, und Brady bemerkt dazu, daß es unmöglich scheine, unter den großen Opern auch nur Eine zu finden, die ein anständiges oder
\quellseite{3}
sauberes Thema habe. In der Regel werde die verbotene Liebe besungen und der Text könne nicht schlechter sein, als er ist. Außerdem seien in jedem dieser elf Stücke noch Lügen, Schwuren, Abenteurer aufgetreten, und Ehescheidungen seien so zahlreich gewesen wie in der heutigen Gesellschaft. Jedes dieser Stücke habe einen schlechten Geschmack im Munde zurückgelassen. Und statt daß nun alle diese Stücke „gepredigt“ und etwa gezeigt hätten, wie die Uebelthäter sich gebessert hätten, habe nur in Einem Stücke der Missethäter ein anderes Leben begonnen, und dieses Stück sei sonst gerade das allerschlechteste und zu Neunzehnteln geradezu viehisch gewesen.

Endlich seien noch zwei Stücke gewesen, die nicht so schlimm waren wie die elf, aber doch auch Anstößiges enthielten. So urtheilt Brady über die besten Theater New Yorks. Die Dutzende schlechter Theater läßt er, wie er selbst sagt und schon bemerkt ist, ganz außer Betracht.

Der andere Mann, der kürzlich ein vernichtendes Urtheil über das heutige Theater abgegeben hat, ist der Redner der St. Louis'er Ethischen Gesellschaft, Walter L. Sheldon. Diese Gesellschaft ist keine religiöse Vereinigung, sie hat kein Glaubensbekenntniß, sie fragt nichts darnach, was ihre Glieder glauben oder nicht glauben, sondern will nur ein „sittliches Leben“ erzielen. Sheldon redet darum lediglich vom Standpunkt eines Weltmannes aus, hält es aber als solcher für nötig, gegen den „geradezu entsetzenerregenden Realismus“ zu protestiren, der heutzutage im Theater dargestellt wird. Mit „Realismus“ meint er die ungeschminkte Darstellung der Sünde und des Lasters und er sagt mit Recht, daß das Gefühl der Einwohner unseres Landes abgestumpft ist, so daß sie nicht mehr erkennen, wie grob die Unzucht in Wort und Handlung im Theater dargestellt wird. „Ich kann nicht verstehen“, bemerkt er, „wie ein junges Mädchen neben einem jungen Mann sitzen kann, ohne vor Scham zu beben, wenn sie manche Worte in den berühmten Schauspielern, wie sie hier in St. Louis aufgeführt werden, hört. Und der Mann, der sie dahin gebracht hat, wird sicherlich erröthen, wenn er sich vergegenwärtigt, in welche Lage er unabsichtlich das junge Mädchen neben sich gebracht hat.“ Sheldon wendet sich dann an die Mütter und fragt sie, ob sie ihre Töchter solchen Erfahrungen aussetzen wollen, zerstört auch den Wahn, als ob solche Schaustücke durch Darstellung der Laster „Moral predigten“ und die Zuschauer von solchen Lastern abhalten würden. So verurtheilt selbst die Welt das heutige Theater, und gerade die hervorragendsten, beliebtesten und am meisten aufgeführten Stücke. Was mögen erst für Stücke und Scenen in den anerkanntermaßen schlechten Theatern zu hören und zu sehen sein! Und was anderes kann die „erzieherische, bildende“ Wirkung des heutigen Theaters sein, als daß es die im Herzen wohnende böse Lust weckt und zu Sünden und Schanden reizt und verführt?

Und dieses heutige Theater sollten nun unsere jungen Christen, unsere Jünglinge und Jungfrauen, besuchen können, ohne Schaden an ihrer Seele zu nehmen, ohne ihr Gewissen zu beflecken, sie, die schon in ihrer Taufe entsagt haben dem Teufel und allen seinen Werken und allem seinem Wesen? Das beantworte dir selbst, lieber junger Christ, und du christliches Elternpaar, das du zum Wächter über deine heranwachsenden Söhne und Töchter gestellt bist und sie darum gerade auch vor dem heutigen Theater warnen und bewahren sollst. L. F.
